\section{Introduction}\label{sec:introduction}

Let $(X,B)$ be a normal projective log canonical pair over an algebraically
closed field $k$ of characteristic zero, and put
$D=K_X+B$. When $\kappa(X,D)\geq0$, the sections of multiples of $D$
determine the Iitaka fibration. The effective problem asks whether one
degree, depending only on the dimension and the allowed boundary
coefficients, already determines that fibration. A uniform degree must
control both the birational geometry of its base and the information
lost in passing between different pluricanonical degrees.

We formulate the conclusion directly as an equality of subfields of
the function field. For an integer $\ell>0$, write
\[
             V_\ell(D)=H^0\bigl(X,\OO_X(\floor{\ell D})\bigr).
\]
When $V_\ell(D)\neq0$, the ratios of its nonzero sections generate a
field $F_\ell(D)\subset k(X)$. Define
\[
 K(D)=k\left(\frac{s}{t}:
       s,t\in V_\ell(D)\setminus\{0\},\
       \ell>0,\ V_\ell(D)\neq0\right).
\]
Thus $F_\ell(D)$ is the function field of the image of the complete
degree-$\ell$ system, viewed inside $k(X)$, and $K(D)$ collects these
fields over all degrees. The reflexive sheaf
$\OO_X(\floor{\ell D})$ makes this definition meaningful even when
$\ell D$ is not Cartier.

\begin{theorem}[Uniform log Iitaka degree]\label{thm:main}
For every integer $d\geq5$ and every finite set
$\Phi\subset[0,1]\cap\Q$, there is an integer $m=m(d,\Phi)>0$
with the following property. Let $k$ be any algebraically closed field
of characteristic zero, and let $(X,B)$ be a normal integral projective
lc $d$-fold pair over $k$. Assume that $B\geq0$, every nonzero coefficient
of $B$ belongs to $\Phi$, $D=K_X+B$ is rational Cartier, and
$\kappa(X,D)\geq0$. Then
\[
                       V_m(D)\neq0,\qquad F_m(D)=K(D)
                       \quad\text{inside }k(X).
\]
\end{theorem}

The bound is existential. The proof uses the good-model theorem of
\cite{LA}, the normal log canonical index theorem of \cite{U}, and
the arithmetic Stein-degree theorem of \cite{SD}. Their precise
statements are recalled in Section~\ref{sec:preliminaries}.
The contribution developed here is the passage from these inputs to
uniform denominators for log fibrations with horizontal boundary,
followed by the exact comparison of rounded section fields.
We state the higher-dimensional range $d\geq5$ to complement the
lower-dimensional results discussed below; the intermediate arguments
are formulated in all fibre dimensions.

\subsection{History and the denominator problem}

Iitaka's work established the asymptotic fibration associated with
pluricanonical systems \cite{Iitaka}. Uniform effectivity asks for a
degree independent of the particular variety. In the general-type
case this became uniform birationality of pluricanonical maps,
proved in arbitrary dimension by Hacon--McKernan, Takayama, and
Tsuji \cite{HaconMcKernan,Takayama,Tsuji}.
Hacon--McKernan--Xu extended effective birationality to big log
canonical adjoints with coefficients in a DCC set
\cite[Theorem~1.3]{HMX}. That theorem includes round-down systems,
the convention used here.

When the Kodaira dimension is smaller than the dimension, the
induced log adjoint on the generic fibre has Kodaira dimension zero.
The canonical bundle formula
transfers the adjoint to the base, where the variation of the fibres
contributes a moduli divisor. Fujino--Mori developed this approach
to effective Iitaka fibrations \cite{FujinoMori}; Viehweg--Zhang
obtained effective results with a surface base and controlled fibre
invariants \cite{ViehwegZhang}. Birkar--Zhang's general effectivity
theorem depends on the dimension, the least nonvanishing pluricanonical
degree of the general fibre, and the middle Betti number of a smooth
model of its canonical cover
\cite[Theorem~1.2]{BZ}. Their effective birationality theorem for
polarized pairs \cite[Theorem~1.3]{BZ} is also the final birational
input in the present proof.

The logarithmic problem introduces horizontal boundary on the
generic fibre. Earlier results include the low-dimensional klt
results of Todorov and Todorov--Xu \cite{Todorov,TodorovXu}, and the
boundedness and birationality results of Hacon--Xu for klt pairs whose
boundary is big over the generic point of the Iitaka base
\cite{HaconXu}. Chen--Han--Liu formulate the effective log
Iitaka problem for lc pairs with DCC coefficients and relate it to
good minimal models and complements
\cite[Conjecture~1.2 and Theorem~1.4]{ChenHanLiu}.
They obtain the dimension-at-most-three case in
\cite[Corollary~1.5]{ChenHanLiu}. Our coefficient set is finite,
and the higher-dimensional argument here uses the three companion
theorems stated above.

A bounded index of the generic fibre does not by itself give the
needed Cartier denominator for the moduli divisor. After a curve
base change one can obtain semistable reduction, but its ramification
degree is generally uncontrolled. The relevant question is whether
the action of inertia on a suitably normalized logarithmic volume
form has bounded order. Only this character needs to be bounded;
the covering degree and the order of the entire automorphism need not be.

Fujino--Mori already control finite characters through the cohomology
of the fibre \cite[Theorem~3.1 and Sections~3.6--3.8]{FujinoMori}.
The companion \cite{U} develops the character method for
boundary-free factors using the singular Beauville--Bogomolov
decomposition. The curve and effective-system arguments of \cite{U4}
treat log Calabi--Yau fibrations in total dimension at most four,
using the index of the whole reduced slc fibre.
We adapt these methods to horizontal boundary in arbitrary fibre
dimension. The geometric structure theorem of Matsumura--Wang
\cite[Theorem~1.3]{MW} supplies a rationally connected factor carrying
the boundary, together with abelian and nonabelian
Beauville--Bogomolov factors. The additional work is to retain the
factor decomposition under semilinear actions and to control the
logarithmic character of the rationally connected factor.

\subsection{The main intermediate result and the proof}

The reusable intermediate statement is a bound for weights over
curves. Let $C$ be a smooth projective complex curve, put $K=\C(C)$,
let $z$ be a uniformizer at $c\in C$, and
let $(V,\Delta)$ be a geometrically integral normal projective lc
$s$-fold pair over $K$, with $\Delta\geq0$. For a positive integer $a$,
a degree-$a$ rational
log trivialization is a nonzero rational $a$-canonical form $\phi$
satisfying $\Div(\phi)+a\Delta=0$.
Its weight at $c$ is the infimum
\[
 w_z(\phi)=\inf_E
 \frac{\ord_E\bigl(\phi\wedge(dz/z)^{\otimes a}\bigr)+a}
      {a\,\ord_E z},
\]
where $E$ ranges over divisorial valuations of the total function
field centred over $c$. These are ordinary canonical orders;
the additive $a$ records the logarithmic pole along the base.
Theorem~\ref{thm:weight} proves that an integer depending only on
$s$ and $a$ clears this weight. Once a degree-$a$ trivialization is
given, no further coefficient parameter is needed.

There are two parts to the weight argument. If a dlt model has a
horizontal coefficient-one component, taking a rational pluriresidue
reduces the fibre dimension. Its Stein factorization may change the
curve field. The arithmetic bound of \cite{SD} controls that change,
so residue gives an induction with bounded denominators.

In the klt case the product theorem of \cite{MW} applies, but the
product cover need not carry the inertia action. We pass to a Galois
comparison cover and prove that a bounded power acts separately on
suitable finite covers of the factors. The proof uses the algebra of
endomorphisms of the tangent sheaf of the comparison cover that preserve
its subsheaf of vector fields tangent to the boundary. The required
vanishing of global vector fields tangent to the boundary on the
rationally connected factor follows from a finite invariant log-volume
measure.

After equivariant semistable reduction, each factor form can be
normalized to have weight zero. Their product has weight zero as
well. If the original form differs from this product by a base function $v$
on a curve cover, and $c'$ is a point over $c$ of ramification index $e$, then
\[
                         e\,w_z(\phi)=\frac{\ord_{c'}(v)}{a}.
\]
The inertia identity forces $e$ to divide a bounded multiple of
$\ord_{c'}(v)$ once the factor characters have bounded order. This
cancels the uncontrolled ramification and bounds the original weight
denominator.

For abelian factors the character acts on a bounded-rank integral
cohomology group. For the other factors we instead use a dlt special
fibre on a good model. Du Bois base change determines its
structure-sheaf cohomology, and coherent Lefschetz gives a fixed
point of a bounded power of the action. A bounded further power
preserves a normal component. The residue there is a log trivialization
in the original degree, so its different has coefficients in a fixed
finite set. The normal lc index theorem of \cite{U}, applied to the
cyclic quotient of this component, bounds the character. This step
works on one normal component and is the reason an index theorem
for a reducible special fibre is unnecessary.

To finish the proof over $\C$, let $f:Y\to Z$ be the semiample contraction of a good model
$(Y,B_Y)$ of $(X,B)$. The normal index theorem gives a bounded positive
integer $p_0$ and a rational function $\psi\in\C(Y)^*$ such that the
canonical bundle formula has the exact presentation
\[
 K_Y+B_Y+\frac1{p_0}\Div(\psi)=f^*(K_Z+B_Z+M_Z).
\]
Here $B_Z$ is the discriminant defined by log canonical thresholds,
and the equality specifies the moduli representative $M_Z$.
Fixing the degree $p_0$ is essential: arbitrary rational principal
shifts can introduce new denominators. Slicing the base by a curve
identifies a coefficient of this moduli divisor with the curve weight,
up to an integer. This proves Theorem~\ref{thm:moduli-denominator}:
a uniform multiple of the moduli trace on a suitable smooth birational
model of $Z$ is nef Cartier.

Birkar--Zhang's theorem then gives a birational system on the base.
The remaining step is to identify complete rounded section spaces
upstairs with those downstairs. This yields equality with $K(D)$
inside the original function field, including ratios initially
appearing in other degrees, and permits descent from $\C$ to any
algebraically closed characteristic-zero field.

Section~\ref{sec:mmp} records the relative minimal-model construction.
Section~\ref{sec:weights} proves the residue reduction for weights.
Sections~\ref{sec:blocks} and~\ref{sec:degeneration} establish the
klt weight bound. Section~\ref{sec:fibration} converts that bound
into moduli denominators, and Section~\ref{sec:completion} proves
Theorem~\ref{thm:main}.
